\documentclass[11pt,a4paper]{article}

% ---------------------------------------------------------------------------
% Ledger-Free Trustless Battleship — research paper
% Compile: pdflatex zk-battleship.tex  (run twice for references)
% Self-contained: standard packages only, embedded bibliography.
% ---------------------------------------------------------------------------

\usepackage[T1]{fontenc}
\usepackage[utf8]{inputenc}
\usepackage{lmodern}
\usepackage[a4paper,margin=2.6cm]{geometry}
\usepackage{amsmath,amssymb,amsthm}
\usepackage{booktabs}
\usepackage{xcolor}
\usepackage{listings}
\usepackage{enumitem}
\usepackage{microtype}
\usepackage[colorlinks=true,linkcolor=black,urlcolor=blue!50!black,citecolor=black]{hyperref}

\newtheorem{claim}{Claim}
\newtheorem{definition}{Definition}

\definecolor{codebg}{gray}{0.96}
\lstdefinestyle{box}{
  backgroundcolor=\color{codebg},
  basicstyle=\ttfamily\small,
  breaklines=true,
  showstringspaces=false,
  frame=single,
  framerule=0pt,
  xleftmargin=8pt,
  framexleftmargin=8pt,
  aboveskip=8pt,belowskip=8pt
}
\lstset{style=box}

\newcommand{\code}[1]{\texttt{\small #1}}
\newcommand{\C}{\ensuremath{C}}

\title{\textbf{Ledger-Free Trustless Battleship:}\\[3pt]
\large Self-Contained Zero-Knowledge Shot Proofs with Public-History Binding}
\author{Hadi Ghazi\\ \small \href{mailto:hadigghazi@gmail.com}{hadigghazi@gmail.com}}
\date{July 2026}

\begin{document}
\maketitle

\begin{abstract}
\noindent
Zero-knowledge proofs let the players of an imperfect-information game answer
questions about hidden state, such as whether a shot in Battleship is a hit or
a miss, without revealing the state and without the ability to answer
dishonestly. The zero-knowledge Battleship implementations we are aware of all
keep the public game state, meaning the shot history, the turn order and the
win condition, on a blockchain, and inherit its latency, its cost and its
reliance on external infrastructure. This paper shows that the straightforward
attempt to drop the ledger is unsound. When the party producing a proof also
supplies the shot history against which ``sunk'' and ``you win'' are evaluated,
a losing player can misreport sunk ships and defer a provable loss for as long
as it likes while still producing proofs that verify. The correction is to bind
the public shot history into the public output of every shot proof. The
verifier fired exactly the shots in that history and can compare it against its
own record, so in a two-player turn-based game no consensus layer is needed to
fix the canonical public state. Sunk detection and the win condition are then
computed inside the proof from the bound history, which makes a whole match,
including the moment it ends, independently checkable from its transcript. We
implement the protocol on the RISC~Zero zkVM as an open-source peer-to-peer
system in which each player runs a local prover and browser interface and the
two are joined by an untrusted message relay. The first move is settled by a
Blum coin flip, mock proving is removed at compile time, and the in-circuit
game logic runs with a control flow that does not depend on the secret board.
On a commodity laptop CPU a shot proof takes between 30 and 47\,s to generate
and about 30\,ms to verify, which is fast enough for turn-based play without
special hardware.
\end{abstract}

\section{Introduction}
\label{sec:intro}

Battleship is a game played on trust. Each player secretly arranges a fleet on
a $10\times10$ grid, the opponent calls out coordinates, and the player holding
the board reports each shot as a hit or a miss. Since only that player can see
the board, nothing stops them from moving a ship out of the way during the game
or refusing to admit that a ship has sunk. This is acceptable among friends but
offers no protection against an opponent who is willing to cheat.

Cryptography removes the need to trust the opponent. Before play begins a player
can \emph{commit} to a board, and can then answer each shot with a
\emph{zero-knowledge proof} that the answer is consistent with the committed
board while disclosing nothing else about it. The commit-and-prove approach is
by now standard, and Battleship has been built on it several times: with
hand-written arithmetic circuits checked by an Ethereum contract
\cite{battlezips,mit6857}, with a zkVM that posts proofs to the NEAR blockchain
\cite{risc0battleship}, and as sample applications on privacy-oriented chains
\cite{aleobattleship,minabattleship}. The same idea of hiding state behind
commitments and proving that actions respect it underlies the on-chain game
Dark Forest \cite{darkforest}.

These systems share one design choice. The public part of the game state, which
cells have been fired on, whose turn it is, how many hits have landed and hence
who has won, is kept on a blockchain, and a smart contract adjudicates
everything that is not secret. The chain adds no secrecy, since the commitments
and proofs already provide that. What it adds is an authoritative public
history. This role is easy to miss, and a protocol that removes the chain
without supplying a substitute fails for a reason that has nothing to do with
the zero-knowledge machinery (Section~\ref{sec:problem}). The predicates ``this
shot sank a ship'' and ``I have lost'' both depend on the shot history, so a
prover free to choose that history can make them come out however it likes. A
losing player can keep producing valid proofs indefinitely and never be shown
to have lost. We call this the \emph{immortal fleet}.

This paper presents and evaluates a Battleship protocol that uses no ledger, no
referee and no trusted third party, and in which the whole game, its ending
included, can be verified. It rests on the following observation, which as far
as we know has not been used before. In a two-party turn-based game, the public
state a verifier needs in order to check a proof about the opponent's hidden
board is just the sequence of shots the verifier itself has fired, and the
verifier already holds that sequence. Binding it into the public output of each
proof therefore does the work that the ledger did.

Concretely, each shot proof takes as input the full public history of shots
fired at the responding player, recomputes the commitment from the private
board so that the player is tied to one board for the whole game, evaluates the
hit, sunk and defeated predicates from the board together with the history, and
writes the history into the proof's public output. The verifier accepts the
proof only when the history it contains matches, in order and element by
element, the shots the verifier actually fired. Every receipt is then
self-contained: given the two commitment receipts and the ordered shot
receipts, an outside party can re-verify a complete match and determine the
winner with nothing else to go on, because the history fields link the receipts
to one another (Section~\ref{sec:audit}).

\paragraph{Contributions.}
\begin{enumerate}[noitemsep,topsep=2pt]
  \item We describe the \emph{public-history binding problem} for ledger-free
        commit-and-prove games and give explicit attacks against the naive
        chain-free protocol, covering forged sunk reports and unbounded
        postponement of defeat (Section~\ref{sec:problem}).
  \item We give a complete two-party protocol, consisting of a board commitment
        with an in-proof legality check, a Blum commit-reveal coin flip
        \cite{blum} for the first move, history-bound shot proofs and an
        in-proof win condition, and we analyse its security properties
        (Sections~\ref{sec:protocol}--\ref{sec:security}).
  \item We implement the protocol on the RISC~Zero zkVM \cite{risc0} with a
        deliberately small trust footprint: each player runs a local prover and
        browser UI, an untrusted stateless relay forwards only public messages,
        mock proving is compiled out, and every routine that reads the secret
        board runs a fixed-shape instruction sequence so that proving time
        reveals nothing about the board (Section~\ref{sec:impl}). The code is
        open source.
  \item We report measurements on a commodity laptop CPU: proving and
        verification times, receipt sizes for both composite and succinct
        receipts, the dependence on history length, and an empirical test of
        witness-independence (Section~\ref{sec:eval}).
\end{enumerate}

We are explicit about what is not new here. The commit-and-prove structure for
Battleship predates this work. The contribution is to remove the ledger while
preserving the integrity guarantees and, through proven termination, slightly
strengthening them, and to provide a complete implementation with measurements.
Zero-knowledge proofs cannot compel a player to keep playing, so aborts remain
possible exactly as in the earlier systems; Section~\ref{sec:limits} treats
this boundary.

\section{Background}
\label{sec:background}

\paragraph{Commitments.}
A commitment scheme lets a party fix a value $v$ by publishing $c =
\mathsf{Com}(v; r)$ that reveals nothing about $v$ (\emph{hiding}) and cannot
later be opened to any $v' \neq v$ (\emph{binding}). We use the standard
salted-hash commitment $\C = \mathrm{SHA\mbox{-}256}(\mathsf{encode}(B) \,\|\, s)$
over a canonical board encoding and a fresh 256-bit salt $s$. Binding rests on
collision resistance and hiding on the unpredictability of the salted preimage,
which is the usual random-oracle-style assumption for this construction.

\paragraph{Zero-knowledge proofs and zkVMs.}
A zero-knowledge proof \cite{gmr} convinces a verifier that a statement holds
without revealing anything beyond its truth, and modern succinct
non-interactive systems \cite{groth16,starks} make this practical for real
applications. A \emph{zkVM} lifts proving from circuits to ordinary programs:
the prover runs a program on private inputs and produces a proof that it
executed correctly. RISC~Zero \cite{risc0} proves the execution of RISC-V
binaries using a STARK \cite{starks}. A proving run produces a \emph{receipt}
made of a \emph{journal}, which holds the public outputs the program chose to
commit, and a \emph{seal}, which is the proof itself. Verification checks the
seal against an \emph{image ID}, a digest of the exact guest binary, so a
verifier learns that this particular program ran and produced this particular
journal, and learns nothing further. SP1 \cite{sp1} and Jolt \cite{jolt} are
comparable systems; \cite{zkpsurvey,zkvmcompilers} survey and benchmark the
landscape.

\paragraph{Battleship.}
Two players each place one ship of every length in $\{5,4,3,3,2\}$, seventeen
cells in all, on a $10\times10$ grid. Ships are straight, axis-aligned, within
bounds and non-overlapping, and touching is permitted. Players take turns
naming a cell of the opponent's grid and learn whether it is a hit or a miss. A
ship is announced sunk once all of its cells have been hit, and a player who
has had all seventeen ship cells hit loses. We add the common rule that a cell
may be fired on at most once per board. This caps a game at 100 shots per side
and, in our protocol, keeps the two players' history sequences aligned.

\section{Related work}
\label{sec:related}

Cryptographic treatments of games are older than practical zero-knowledge
proofs. Mental poker \cite{sra81} dealt a deck between mutually distrusting
players, and the completeness theorems for secure computation \cite{gmw87} show
that in principle any game can be played fairly. What has changed since is
practicality: succinct proofs and zkVMs have made per-move proving cheap enough
to use.

\emph{On-chain ZK Battleship.} BattleZips \cite{battlezips}, built on circom
and Groth16, proves board legality at placement and hit-or-miss on each shot,
while an Ethereum contract enforces turn order, records declared shots, counts
the seventeen hits that end the game, and is the root of trust for all public
state. The MIT 6.857 project \cite{mit6857}, using circom and snarkjs on
Ethereum, proves that each reported result is consistent with the hashed board
but, unlike BattleZips, does not prove that the committed board is legal, so a
player may commit a malformed board; the win condition is again computed by the
contract. RISC~Zero's own example \cite{risc0battleship} put the logic in a
zkVM but stored game state in a NEAR contract, and is now unmaintained and
built against an old receipt format. Aleo and Mina host comparable
demonstrations \cite{aleobattleship,minabattleship}. In each case the chain is
what pins the public history, and this paper is about what takes its place once
the chain is gone.

\emph{Off-chain ZK gaming.} Dark Forest \cite{darkforest} was the first
deployed game to hide state behind commitments, though its core still runs
on-chain. State-channel constructions \cite{statechannels,hawk} carry out
interaction off-chain between on-chain checkpoints, and ``ZK state channels''
have been proposed for turn-based games, but they still rely on a chain to
settle disputes. Our setting is more restrictive, in that no chain is present
at any layer, and the point we make is that for two-player turn-based games a
chain is unnecessary for integrity and is useful only as a remedy against
aborts (Section~\ref{sec:limits}).

\emph{Fair coin flipping.} The first move is chosen with Blum's commit-reveal
protocol \cite{blum}. Its familiar weakness, that the party revealing last can
abort after seeing an unfavourable result, is already covered by the fact that
any zero-knowledge game protocol permits aborts.

\section{The public-history binding problem}
\label{sec:problem}

Consider the natural ledger-free protocol, which up to notation is what the
systems above become once the contract is removed. Player $P$ holds a board $B$
and salt $s$ and has published the commitment $\C = H(\mathsf{encode}(B)\|s)$.

\begin{itemize}[noitemsep,topsep=2pt]
\item \textbf{Setup proof.} $P$ proves knowledge of a \emph{legal} $B$ and
      $s$ with $H(\mathsf{encode}(B)\|s) = \C$; the journal is $\C$.
\item \textbf{Naive shot proof.} On shot $x$, $P$ proves: ``for the $B,s$
      behind $\C$: $\mathit{hit} = \mathsf{occ}(B,x)$, and
      $\mathit{sunk} = \mathsf{sunk}(B, x, \mathcal{H})$'', where
      $\mathcal{H}$ is the set of shots previously fired at $P$, supplied by
      $P$ as a private input. The journal is
      $(\C, x, \mathit{hit}, \mathit{sunk})$.
\end{itemize}

The hit bit is sound, since it is a function of $B$ and $x$ alone and both are
bound by the journal. Nothing that depends on $\mathcal{H}$ is sound, because
the verifier cannot tell which $\mathcal{H}$ the statement was evaluated
against. This gives three attacks.

\paragraph{A1: Sunk suppression.} On the shot that would complete a ship, a
dishonest responder proves against a history $\mathcal{H}' \subsetneq
\mathcal{H}$ that omits an earlier hit. The proof verifies and reports
$\mathit{hit}=\textsf{true}$, $\mathit{sunk}=\bot$, so the true state of the
responder's fleet is hidden. Where sunk announcements affect strategy this is a
real advantage, and it costs the attacker nothing.

\paragraph{A2: Phantom sunk.} In the other direction, the responder adds to
$\mathcal{H}'$ shots that were never fired, so that a ship is announced sunk
before it should be. This can be used to feign weakness or to disrupt an
opponent who is tracking sunk counts.

\paragraph{A3: The immortal fleet.} The win condition is
$\mathsf{allsunk}(B, \mathcal{H})$, and whoever evaluates it needs
$\mathcal{H}$. Without a ledger, an implementation must either believe the
loser's client when it declares defeat or have the winner's client tally sunk
reports, and A1 already undermines the latter. A losing player is therefore
never \emph{provably} beaten: it can answer every shot with a valid proof,
withhold the last sunk reports, and the game simply does not end. No
cryptographic assumption fails; the proofs are true statements about a history
that was never agreed upon. An earlier version of our own implementation had
precisely this defect, with defeat recorded as a self-reported boolean outside
any proof, and that is what led us to the present analysis.

On-chain systems avoid A1 through A3 because the contract is the history: shots
are transactions and hits are counted in consensus. The problem is to find what
plays that role when no chain is present.

\paragraph{The observation.} In a two-player turn-based game the history
$\mathcal{H}$ that matters for proofs about $P$'s board is exactly the sequence
of shots fired by the opponent $V$, and $V$ is the party verifying those proofs
and has chosen and recorded every one of them. The canonical public state
already sits with the party that needs it, so it calls for binding rather than
consensus. The fix follows: move $\mathcal{H}$ from the witness into the
statement by writing it into the journal, and have $V$ check it against its own
record.

\section{The protocol}
\label{sec:protocol}

We now describe the protocol as implemented. $H$ is SHA-256, and
$\mathsf{encode}(B)$ lists the five ships sorted by a canonical ship
identifier, each as (id, row, col, orientation) in 4 bytes, followed by the
32-byte salt, so that the commitment does not depend on the order in which
placements happen to be listed.

\subsection{Setup}

Each player $P \in \{A, B\}$:
\begin{enumerate}[noitemsep,topsep=2pt]
\item chooses a board $B_P$ (interactively, in the UI) and samples
      $s_P \leftarrow \{0,1\}^{256}$ from the OS CSPRNG;
\item runs the \textsf{commit} guest, which \emph{asserts}
      $\mathsf{legal}(B_P)$ (composition, bounds, non-overlap) and outputs
      journal $\C_P = H(\mathsf{encode}(B_P)\|s_P)$, yielding receipt
      $R^{\mathsf{com}}_P$;
\item samples a coin nonce $n_P \leftarrow \{0,1\}^{256}$ and sends
      $(\C_P, R^{\mathsf{com}}_P, H(n_P))$ to the opponent;
\item on receiving the opponent's tuple, verifies $R^{\mathsf{com}}$ against
      the pinned \textsf{commit} image ID, checks that the journal equals the
      claimed $\C$, and rejects a mirrored commitment $\C = \C_P$ (a copied
      commitment can never be answered, since the copier does not hold the
      witness, so a game against it could only stall);
\item reveals $n_P$; verifies the opponent's reveal against $H(n)$;
      the first mover is decided by the XOR parity of the two nonces
      \cite{blum}.
\end{enumerate}

An illegal board therefore never enters a game, because no receipt exists for
one. This is in contrast to \cite{mit6857}, where placement legality is not
proved.

\subsection{Play}

Let $F_V = (x_1, \dots, x_k)$ be the ordered shots $V$ has fired at $P$ so
far (the \emph{fired list}), maintained by $V$'s own agent. Shots must be
pairwise distinct, and $P$'s agent refuses to answer a repeated cell.

On $V$'s turn, $V$ sends a shot $x$. $P$ runs the \textsf{shot} guest on
private input $(B_P, s_P)$ and public input $(\C_P, x, \mathcal{H})$, where
$\mathcal{H} = (x_1, \dots, x_k)$ is $P$'s record of shots received. The
guest:
\begin{enumerate}[noitemsep,topsep=2pt]
\item recomputes $H(\mathsf{encode}(B_P)\|s_P)$ and \emph{asserts} it equals
      $\C_P$ \hfill (board binding)
\item computes $\mathit{hit} = \mathsf{occ}(B_P, x)$;\quad
      $\mathit{sunk} = \mathsf{sunkship}(B_P, x, \mathcal{H})$;\quad
      $\mathit{def} = \mathsf{allsunk}(B_P, \mathcal{H} \cup \{x\})$
\item commits journal
      $J = (\C_P,\, x,\, \mathit{hit},\, \mathit{sunk},\, \mathit{def},\,
      \mathcal{H})$ \hfill (history binding)
\end{enumerate}
and returns the receipt. $V$ accepts iff:
\[
\mathsf{Verify}(R, \mathsf{id}_{\mathsf{shot}}) \;\wedge\;
J.\C = \C_P \;\wedge\; J.x = x \;\wedge\; J.\mathcal{H} = F_V ,
\]
where the last conjunct is \emph{sequence} equality against $V$'s own fired
list, and $V$ appends $x$ to $F_V$ only after accepting. When $J.\mathit{def} =
\textsf{true}$, $V$ has won and holds a proof of the win that it can show to
anyone; $P$'s own agent reads the same bit from the same journal, so the defeat
is at once undeniable and impossible to fake. The roles then swap.

Since the responder's history record and the verifier's fired list each grow by
exactly the shot that was accepted, they remain equal by induction. Any
divergence, whether a dropped, injected, reordered or replayed message, and
whether caused by the network, the relay or a cheating peer, causes the next
verification to fail.

\subsection{Public auditability of transcripts}
\label{sec:audit}

A match transcript consists of the two commitment receipts together with, for
each direction of play, the ordered sequence of shot receipts. An uninvolved
third party can check, using only the pinned image IDs, that (i) every receipt
verifies, (ii) within a direction the journal history of receipt $k$ equals the
shots of receipts $1$ through $k-1$, so the receipts chain to one another with
no need for signatures or timestamps, and (iii) the final receipt of the losing
direction carries $\mathit{def} = \textsf{true}$. The end of the game is thus a
property of the transcript rather than an assertion by a player. This is the
ledger-free counterpart of a contract declaring the game over, and it is
exactly what the naive protocol cannot provide.

\subsection{Cost}

Each shot costs one proof, one verification, and a journal holding
$O(|\mathcal{H}|)$ shot entries at two bytes each before serialization. Across a
full 100-shot direction the journal overhead summed over the transcript is
$O(n^2)$, on the order of $10^4$ bytes, which is small next to the seals
themselves (Section~\ref{sec:eval}). A running-hash binding
($h_k = H(h_{k-1} \| x_k)$) would reduce each journal to $O(1)$ but force an
auditor to recompute the chain; for $n \le 100$ the explicit list is simpler
and self-describing.

\section{Security analysis}
\label{sec:security}

The adversary is a computationally bounded player who may deviate however it
wishes: it can send any messages, commit to any board, present any history, and
abort at any time. The relay is trusted only to deliver messages; it can
observe and drop them but never carries anything other than public data. We
assume the collision resistance of SHA-256, the soundness and zero-knowledge of
the RISC~Zero STARK \cite{risc0,starks}, and the correctness of the roughly
300-line shared rules crate, which is compiled into the guest and constitutes
the protocol's entire trusted rulebook; both players pin its image IDs and
check them on every verification.

\begin{claim}[Response soundness]
If $V$ accepts journal $J$ for shot $x$, then except with negligible
probability there exist $B, s$ with $H(\mathsf{encode}(B)\|s) = \C_P$,
$\mathsf{legal}(B)$, and $J.\mathit{hit}, J.\mathit{sunk}, J.\mathit{def}$
equal to the rule functions evaluated at $(B, x, F_V)$.
\end{claim}
\emph{Sketch.} STARK soundness gives an execution of the pinned guest on some
witness $(B,s)$ that produces $J$. The guest's first assertion forces
$H(\mathsf{encode}(B)\|s) = J.\C = \C_P$, and legality of $B$ follows from the
setup receipt for $\C_P$ together with collision resistance, since two distinct
boards behind one $\C$ would be a collision. The guest evaluates the rule
functions on $(B, x, J.\mathcal{H})$ by construction, and $V$'s check
$J.\mathcal{H} = F_V$ replaces the journal's history with the true one.
\hfill$\square$

\begin{claim}[Board binding across the match]
All accepted journals for commitment $\C_P$ are consistent with a single
board $B$, except with negligible probability.
\end{claim}
\emph{Sketch.} Each accepted journal is backed by a witness that hashes to
$\C_P$, and two witnesses carrying different boards would form a SHA-256
collision, since the encoding is injective on legal boards: it records each
ship's kind, anchor and orientation in a canonical order. \hfill$\square$

\begin{claim}[Termination soundness and undeniability]
$V$ obtains a transferable proof of victory exactly when
$F_V \cup \{x\}$ covers all 17 ship cells of $P$'s committed board; $P$
cannot postpone this by any choice of messages other than refusing to
answer.
\end{claim}
\emph{Sketch.} $\mathit{def}$ is computed inside the guest from $(B, x,
\mathcal{H})$ with $\mathcal{H}$ bound to $F_V$, so by response soundness it
matches the true predicate. To avoid a journal with $\mathit{def} =
\textsf{true}$, a responder would have to present a different history, which
$V$'s comparison rejects, a different board, which is a collision, or no proof
at all, which is an abort and is visible as one. Attacks A1 through A3 are ruled
out in the same way. \hfill$\square$

\begin{claim}[Secrecy]
A (possibly malicious) verifier interacting with an honest prover learns
nothing about the prover's board beyond the game-inherent information: the
hit/miss/sunk/defeated values of the queried cells, and what those imply.
\end{claim}
\emph{Sketch.} The prover's messages are receipts, and by the zero-knowledge
property of the proof system the verifier's view can be simulated from the
journals alone. Those journals contain only $\C$, which hides the board through
the 256-bit salt, the verifier's own shots, and the game-inherent result bits.
The guarantee is deliberately one of a leaky function: the answers in
Battleship themselves shrink the space of possible boards, which is the game
working as intended, and the protocol protects only what remains. \hfill$\square$

\paragraph{Replay, reflection, reordering.} Journals bind $(\C, x,
\mathcal{H})$, so a replayed receipt fails on the $x$ or $\mathcal{H}$ check,
and a receipt from the opposite direction of play refers to the wrong
commitment. An opponent that mirrors a commitment is rejected during setup and
could not answer shots in any case.

\paragraph{Malicious relay.} The relay carries commitments, coin values, shot
coordinates and receipts, all of which are public. It cannot change any of them
without detection, since verification fails on a mismatch; it cannot forge
outcomes, since it cannot produce receipts for the pinned image IDs; and it
learns nothing that the transcript does not already contain. Its only real
power is to deny service.

\paragraph{Prover-side side channels.} Proof contents reveal nothing beyond the
journals, but the time taken to prove is visible to the opponent. If the
guest's control flow depended on ship positions, this latency could leak
information about the board. Every routine in the shared crate that reads the
board is therefore written with fixed-shape loops, using the fact that a legal
fleet always occupies exactly seventeen cells, and with boolean accumulation
that does not short-circuit, so that the length of the instruction trace does
not depend on the witness. Section~\ref{sec:eval} checks this invariance
empirically. We make no claim about physical side channels on the prover's own
machine, which lie outside our scope.

\section{Implementation}
\label{sec:impl}

The system comprises roughly 3{,}400 lines of Rust and TypeScript and is open
source under the MIT license.

\begin{itemize}[noitemsep,topsep=2pt]
\item \code{core}: the shared rules crate discussed above, defining the types,
      fleet legality, the hit, sunk and defeat predicates, and the canonical
      encoding. It is compiled unchanged into both the native host and the
      RISC-V guest, so the rules that are proved are identical to the rules the
      application plays. Its unit tests include an equivalence suite that checks
      the fixed-shape implementations against straightforward reference versions
      over random boards.
\item \code{methods/guest}: the two guest programs, \textsf{commit} and
      \textsf{shot}, each about 40 lines on top of \code{core}. SHA-256 uses
      the zkVM's accelerated implementation.
\item \code{host}: the prover and verifier library together with the
      \emph{prover agent}, a localhost-only HTTP service that holds the
      player's board and salt, proves the player's answers, and verifies the
      opponent's, checking the receipt, image ID, commitment, shot and history.
      The build turns on RISC~Zero's \code{disable-dev-mode}, and this together
      with clearing \code{RISC0\_DEV\_MODE} from the environment makes mock
      receipts impossible rather than merely discouraged. Adversarial
      integration tests run each attack from Section~\ref{sec:problem} against
      the real prover and check that it fails.
\item \code{relay}: a WebSocket forwarder of about 200 lines, with no game
      logic and no dependency on RISC~Zero, deployable as a container anywhere.
      It pairs two connections by room code and passes messages between them.
\item \code{web}: a Next.js interface for each player, covering fleet
      placement, firing and status, which communicates only with that player's
      own agent and with the relay. The Blum coin flip runs here through
      WebCrypto, within the trust domain of the player it protects.
\end{itemize}

Two properties of the deployment are worth stating. First, a secret is confined
to a single process: the board and salt exist only inside the player's own
agent, which binds to \code{127.0.0.1} and answers CORS requests only from the
local UI origin. Second, each agent exposes its guest image IDs at
\code{/health}, so players can compare them out of band, although a mismatch is
in any event unexploitable and merely causes every cross-verification to fail.

\section{Evaluation}
\label{sec:eval}

\paragraph{Setup.} All measurements are CPU-only on a consumer laptop:
Intel Core Ultra~7 155H (22 hardware threads), 15\,GB RAM available to the
measurement environment (WSL2, Ubuntu 22.04), \code{rustc} 1.96,
\code{risc0-zkvm} 3.0.5, default composite receipts and
\code{ProverOpts::succinct()} where indicated. Wire size is the
JSON-serialized receipt as actually sent through the relay.

\begin{table}[h]
\centering
\small
\begin{tabular}{@{}llrrrrrr@{}}
\toprule
\textbf{Proof} & \textbf{Receipt} & \textbf{$|\mathcal{H}|$} &
\textbf{Prove (s)} & \textbf{Verify (ms)} & \textbf{Journal (B)} &
\textbf{Seal (KiB)} & \textbf{Wire (KiB)} \\
\midrule
commit & composite & --  & 21.0 & 34.5 & 128 & 215.7 & 554.1 \\
shot   & composite & 0   & 30.0 & 34.3 & 152 & 215.7 & 554.1 \\
shot   & composite & 25  & 47.2 & 27.1 & 352 & 238.2 & 613.6 \\
shot   & composite & 50  & 41.1 & 28.2 & 552 & 238.2 & 614.4 \\
shot   & composite & 99  & 42.3 & 32.8 & 944 & 238.2 & 616.1 \\
commit & succinct  & --  & 54.3 & 31.1 & 128 & 217.4 & 569.5 \\
shot   & succinct  & 0   & 54.8 & 29.0 & 152 & 217.4 & 569.5 \\
shot   & succinct  & 99  & 76.8 & 27.1 & 944 & 217.4 & 572.8 \\
\bottomrule
\end{tabular}
\caption{Proving and verification cost (single runs; proving uses all 22
threads). $|\mathcal{H}|$ is the length of the bound shot history. Wire size
is the JSON-serialized receipt as sent through the relay; Seal is the raw
STARK seal within it.}
\label{tab:bench}
\end{table}

\paragraph{Latency and playability.} Setup requires one commit proof per
player, about 21\,s and produced concurrently on the two machines, plus a coin
flip that takes under a millisecond. Each subsequent turn requires one shot
proof, between 30 and 47\,s, along with verification of about 30\,ms and the
relay round-trips. Guest execution grows linearly in the bound history, from
20.3k user cycles at $|\mathcal{H}| = 0$ to 101.8k at $|\mathcal{H}| = 99$, but
the prover pads execution into segments of a power-of-two length, so the cost
of proving changes only when a segment boundary is crossed: an empty history
occupies one segment of $2^{16}$ cycles, and every longer history in a full
game fits within two. Binding the history therefore adds almost nothing to the
proving time, and its journal overhead stays under 1\,KB against a receipt of
about 600\,KiB even at $|\mathcal{H}| = 99$, so the $O(n^2)$ transcript overhead
noted in Section~\ref{sec:protocol} does not matter in practice. A complete
match of roughly 60 to 90 shots spends between 30 and 60 minutes proving, spread
across the game at the pace of a correspondence match, and no single wait
exceeds a minute on this hardware. More cores, a GPU, or a newer prover
\cite{zkvmcompilers} would reduce these times.

Succinct receipts compress an entire execution into a seal of constant size,
217.4\,KiB regardless of the history, at between 1.4 and 1.8 times the proving
cost. At this circuit size a composite receipt spans only one or two segments,
so succinctness saves nothing on the wire, and we use composite receipts by
default. Recursion becomes worthwhile only when many moves are to be aggregated
(Section~\ref{sec:limits}).

\paragraph{Witness independence.} Running the \textsf{shot} guest on a fixed
shot and a 30-shot history across 20 random boards gives execution lengths
between 45{,}147 and 45{,}173 cycles, a spread of 26 cycles, below 0.06\%,
attributable to runtime bookkeeping outside the rule functions. The spread
disappears once proving begins, since all 20 executions pad to the same
$2^{16}$-cycle segment and therefore take the same time to prove, leaving no
signal about the board. With the fixed-shape rule implementations, this gives
empirical support for the side-channel claim of Section~\ref{sec:security}
rather than resting on inspection of the code alone.

\paragraph{Comparison with on-chain designs.} A direct quantitative comparison
with on-chain designs would not be meaningful, given the different trust models
and hardware, but the qualitative difference is clear. On-chain systems pay
consensus latency, on the order of tens of seconds per move on proof-of-work
Ethereum in \cite{mit6857}, together with a transaction fee per move, and they
place game metadata in public permanently. Our design pays only local proving
time and the relay round-trips, keeps the transcript private to the two
players, who may publish it if they want it audited, and forgoes the chain's
built-in penalty for aborting, namely stake slashing, which we turn to next.

\section{Limitations and scope}
\label{sec:limits}

\paragraph{Aborts.} No zero-knowledge protocol can compel a losing player to
keep generating proofs, and a player facing defeat can simply disconnect. What
the protocol does guarantee is that such an abort is attributable and cannot be
faked, since the transcript shows who stopped answering and that every answered
shot was honest, and that the outcome of a completed game is not open to
dispute. Remedies such as forfeit timers by social convention, or stakes with
slashing for anyone who chooses to add a chain or an escrow, sit outside the
protocol and compose with it unchanged.

\paragraph{Trusted rulebook.} Soundness is relative to the correctness of the
roughly 300-line rules crate and of the zkVM toolchain. The crate is small, is
shared verbatim between play and proof, and is tested adversarially; trust in
the toolchain is common to the whole zkVM ecosystem \cite{zkpsurvey}.

\paragraph{Metadata.} The protocol conceals boards rather than behaviour.
Timing, choice of moves and disconnections are visible to the opponent and the
relay, and proving-time side channels are reduced by the fixed-shape traces but
not physically eliminated. Reusing a board or a salt across games invalidates
the hiding argument, and the implementation draws a fresh salt for each game and
advises against reusing boards.

\paragraph{Generality.} The construction extends directly to any two-player
turn-based game with hidden state whose public state is a deterministic function
of the public sequence of moves, such as Stratego-style reveals or fog-of-war
variants, by substituting the rule functions inside the guest. Games with
hidden moves, simultaneous turns, or more than two players bring back the need
for a shared ordering or a fairness layer and are outside our scope.

\section{Conclusion}
\label{sec:conclusion}

Deployed zero-knowledge Battleship systems hand the public half of the game,
its history, turns and termination, to a blockchain. We have shown that removing
the chain without more care breaks exactly that half, producing unverifiable
sunk reports and the immortal fleet, and that for two-player turn-based games no
consensus is needed to repair it. The verifier already holds the canonical
history, and binding it into each proof's public output, together with computing
the win condition inside the proof, yields self-contained receipts, matches that
terminate verifiably, and a deployment that consists of nothing more than the
two players and a message relay. The implementation runs on laptop CPUs and is
open source, and we intend it as a compact and complete example of a trustless
hidden-state game built without a chain.

\subsection*{Artifact}

Source code (protocol, agent, relay, UI, benchmarks, and the LaTeX for this
paper) is available at
\url{https://github.com/hadigghazi/Ledger-Free-Trustless-Battleship}, and the
v1.0.0 release is archived at
\href{https://doi.org/10.5281/zenodo.21207531}{DOI 10.5281/zenodo.21207531}.
All reported numbers are reproducible via
\code{cargo run --release -p host --bin zk-battleship-bench}
and the \code{--ignored} integration test suite.

\begin{thebibliography}{99}\small

\bibitem{gmr} S.~Goldwasser, S.~Micali, and C.~Rackoff.
The knowledge complexity of interactive proof systems.
\emph{SIAM Journal on Computing}, 18(1):186--208, 1989.

\bibitem{sra81} A.~Shamir, R.~L. Rivest, and L.~M. Adleman.
Mental poker. In \emph{The Mathematical Gardner}, pages 37--43. Springer, 1981.

\bibitem{gmw87} O.~Goldreich, S.~Micali, and A.~Wigderson.
How to play any mental game, or a completeness theorem for protocols with
honest majority. In \emph{Proc.\ 19th ACM STOC}, pages 218--229, 1987.

\bibitem{blum} M.~Blum.
Coin flipping by telephone: a protocol for solving impossible problems.
\emph{ACM SIGACT News}, 15(1):23--27, 1983.

\bibitem{starks} E.~Ben-Sasson, I.~Bentov, Y.~Horesh, and M.~Riabzev.
Scalable, transparent, and post-quantum secure computational integrity.
Cryptology ePrint Archive, Report 2018/046, 2018.

\bibitem{groth16} J.~Groth.
On the size of pairing-based non-interactive arguments.
In \emph{EUROCRYPT 2016}, LNCS 9666, pages 305--326. Springer, 2016.

\bibitem{risc0} J.~Bruestle and P.~Gafni.
RISC Zero zkVM: scalable, transparent arguments of RISC-V integrity.
RISC Zero, 2023. \url{https://dev.risczero.com/proof-system-in-detail.pdf}.

\bibitem{sp1} Succinct Labs.
SP1: a performant, open-source zkVM.
2024. \url{https://github.com/succinctlabs/sp1}.

\bibitem{jolt} A.~Arun, S.~Setty, and J.~Thaler.
Jolt: SNARKs for virtual machines via lookups.
In \emph{EUROCRYPT 2024}, LNCS 14656, pages 3--33. Springer, 2024.

\bibitem{zkpsurvey} R.~Lavin, X.~Liu, H.~Mohanty, L.~Norman, G.~Zaarour, and
B.~Krishnamachari.
A survey on the applications of zero-knowledge proofs.
arXiv:2408.00243, 2024.

\bibitem{zkvmcompilers} T.~Gassmann, S.~Chaliasos, T.~Sotiropoulos, and
Z.~Su. Evaluating compiler optimization impacts on zkVM performance.
In \emph{Proc.\ ASPLOS 2026}. ACM, 2026. arXiv:2508.17518.

\bibitem{battlezips} BattleZips.
BattleZips-Circom: zero-knowledge arithmetic circuits for on-chain
Battleship. 2022. \url{https://github.com/BattleZips/BattleZips-Circom}.

\bibitem{mit6857} A.~Gupta, N.~Kaashoek, B.~Wang, and J.~Zhao.
Zero-knowledge Battleship. MIT 6.857 course project report, 2020.
\url{https://courses.csail.mit.edu/6.857/2020/projects/13-Gupta-Kaashoek-Wang-Zhao.pdf}.

\bibitem{risc0battleship} RISC Zero.
Battleship example on NEAR. 2022.
\url{https://github.com/risc0/battleship-example}.

\bibitem{aleobattleship} Aleo.
ZK Battleship in Leo (workshop example). 2023.
\url{https://github.com/AleoNet/workshop}.

\bibitem{minabattleship} Mina Foundation.
Battleships: a zero-knowledge game on Mina with o1js. 2023.
\url{https://github.com/o1-labs}.

\bibitem{darkforest} Dark Forest team.
Announcing Dark Forest. 2020. \url{https://blog.zkga.me/announcing-darkforest}.

\bibitem{hawk} A.~Kosba, A.~Miller, E.~Shi, Z.~Wen, and C.~Papamanthou.
Hawk: the blockchain model of cryptography and privacy-preserving smart
contracts. In \emph{IEEE S\&P 2016}, pages 839--858, 2016.

\bibitem{statechannels} J.~Coleman, L.~Horne, and L.~Xuanji.
Counterfactual: generalized state channels. 2018.
\url{https://l4.ventures/papers/statechannels.pdf}.

\end{thebibliography}

\end{document}
